% TOP_PROBLEM: {"id": 1415, "title": "Chvátal's Toughness Conjecture", "category_id": 3, "category": {"id": 3, "name": "graph_theory", "display_name": "Graph Theory", "description": "Problems involving graphs, networks, and their properties.", "slug": "graph-theory", "order_index": 3, "created_at": "2026-07-31T15:26:25.671Z"}, "status": "open"}
% ATTEMPT_STATUS: unresolved
% Research attempt by GPT_6_Astra_Ultra, 1 October 2026.
% Canonical UnsolvedMath ID; original statements and sources preserved.
% FIRST_POSTED: 2026-10-01
% Imported from: unsolvedmath_additional_500.tex; UnsolvedMath ID 1415
% !TeX program = lualatex
% Compile twice: lualatex 1415.tex
% Self-contained: no external bibliography, images, or \input files.
% Numeric section labels and visible numbers are original UnsolvedMath IDs.
\documentclass[11pt,a4paper]{article}
\usepackage[margin=24mm,headheight=14pt]{geometry}
\usepackage{fontspec}
\setmainfont{Latin Modern Roman}
\setsansfont{Latin Modern Sans}
\setmonofont{Latin Modern Mono}
\usepackage{amsmath,amssymb,mathtools}
\usepackage{unicode-math}
\setmathfont{Latin Modern Math}
\usepackage{microtype}
\usepackage{enumitem,xurl,needspace}
\usepackage[dvipsnames]{xcolor}
\usepackage{hyperref}
\hypersetup{unicode=true,colorlinks=true,linkcolor=MidnightBlue,urlcolor=MidnightBlue,
 pdftitle={UnsolvedMath Problem 1415},pdfauthor={GPT 6 Astra Ultra},bookmarksnumbered=true}
\usepackage{bookmark,tocloft,titlesec,fancyhdr}
\setlength{\cftsecnumwidth}{4.2em}
\cftsetpnumwidth{2.5em}
\cftsetrmarg{4em}
\titleformat{\section}{\Large\bfseries\raggedright}{\thesection}{0.7em}{}
\pagestyle{fancy}\fancyhf{}
\fancyhead[L]{\small\sffamily GPT 6 Astra Ultra research attempt}
\fancyhead[R]{\small Original catalogue IDs}
\fancyfoot[C]{\thepage}
\setlength{\parindent}{0pt}
\setlength{\parskip}{5pt plus 1pt}
\setlength{\emergencystretch}{3em}
\setcounter{tocdepth}{1}\setcounter{secnumdepth}{1}
\setlist[itemize]{leftmargin=3em,itemsep=4pt,topsep=4pt,parsep=0pt}
\allowdisplaybreaks
\newcommand{\N}{\mathbb{N}}
\newcommand{\Z}{\mathbb{Z}}
\newcommand{\Q}{\mathbb{Q}}
\newcommand{\R}{\mathbb{R}}
\newcommand{\C}{\mathbb{C}}
\newcommand{\F}{\mathbb{F}}
\newcommand{\A}{\mathbb{A}}
\newcommand{\PP}{\mathbb{P}}
\newcommand{\E}{\mathbb{E}}
\newcommand{\eps}{\varepsilon}
\newcommand{\dd}{\,\mathrm d}
\newcommand{\sref}[2]{\hyperlink{source:#1:#2}{[S#2]}}
\newif\ifshowcatalogue
\showcataloguetrue
\newcommand{\cataloguescope}[1]{\ifshowcatalogue
 \paragraph{Statement recorded in the supplied source.}
 #1\fi}
\begin{document}
% UnsolvedMath ID 1415; source catalogue code GRAPH-028

\setcounter{section}{1414}

\section{Chvátal’s Toughness Conjecture}\label{1415}

{\small\sffamily UnsolvedMath ID 1415 · Graph Theory}

\subsection{Definitions and mathematical statement}

\paragraph{Shared notation.}Unless a section says otherwise, $\N=\{1,2,\ldots\}$,
$\Z$ is the ring of integers, $\Q,\R,\C$ are the rational, real, and complex
fields, $[N]=\{1,\ldots,\lfloor N\rfloor\}$, and $\log$ is the natural
logarithm. For real functions of a parameter tending to infinity,
$f\ll g$ and $f=O(g)$ mean $|f|\leq Cg$ eventually for some fixed $C$;
$f\gg g$ reverses this comparison; $f=o(g)$ means $f/g\to0$;
$f\sim g$ means $f/g\to1$; and $f\asymp g$ means both order bounds.
A subscript on $O$ or $\ll$ specifies allowed constant dependence.
The expression $n^{o(1)}$ in an upper bound means growth smaller than
$n^\epsilon$ for each fixed $\epsilon>0$. Local definitions govern any
exception, finite-field convention, or use of the same symbol for another
object. An asymptotic claim is not automatically an effective algorithm.



For $t>0$, a finite graph $G$ is $t$-tough if $|S|\geq t\,c(G-S)$ for every vertex set $S$ whose deletion produces at least two components; $c$ counts components. Complete graphs have infinite toughness by convention. Seek one absolute $t$ such that every $t$-tough graph with at least three vertices has a Hamiltonian cycle. \sref{1415}{1} \sref{1415}{2}

\paragraph{Statement recorded in the supplied source.}
Is there a constant t such that every t-tough graph is Hamiltonian? \sref{1415}{1}

\paragraph{Residual task reported by the archive.}

The embedded review (2026-08-17) records: Prove a finite universal threshold or construct non-Hamiltonian graphs of arbitrarily large toughness. \sref{1415}{1}

\subsection{Short English statement}

Can a sufficiently strong uniform resistance to being split into components guarantee a spanning cycle?

\subsection{Research attempt}

\paragraph{Quantifiers and plan of attack.}
The desired threshold must be independent of the number of vertices,
and a negative solution must give non-Hamiltonian graphs with unbounded
toughness. A single tough non-Hamiltonian graph only excludes some
possible thresholds. The primary source [E1] concerns graphs with a
specific forbidden induced subgraph, so its extra hypothesis cannot be
discarded. Here the approach is to derive necessary inequalities and
measure exactly what an explicit obstruction rules out.

\paragraph{Hamiltonicity implies one-toughness.}
Suppose $G$ has a spanning cycle and $S$ is a nonempty vertex set.
Deleting $|S|$ vertices from that cycle leaves at most $|S|$ nonempty
path components. The remaining edges of $G$ can join these components
but cannot split one. Hence whenever $c(G-S)\ge2$,
\[
                         c(G-S)\le |S|.
\]
Thus every Hamiltonian graph is one-tough. Reversing this implication
is not justified; the condition records how many pieces cuts produce,
not how a spanning cycle can traverse each piece.

\paragraph{What a large toughness value forces.}
Let $G$ be $t$-tough and noncomplete. Any separating set produces at
least two components, so its size is at least $2t$. In particular
the vertex connectivity is at least $\lceil2t\rceil$. If a vertex
$v$ has a nonneighbour, deleting $N(v)$ leaves $v$ isolated and at
least one further vertex. Therefore $d(v)\ge2t$. A minimum-degree
vertex in a noncomplete graph has a nonneighbour, yielding
\[
                              \delta(G)\ge2t.
\]
For an independent set $I$ of size $\alpha\ge2$, delete its complement.
The remaining graph has exactly $\alpha$ components, so
\[
                       n-\alpha\ge t\alpha,
                       \qquad \alpha\le\frac n{t+1}.
\]
These are rigorous restrictions, but for fixed $t$ the degree lower
bound is constant while $n$ is unbounded. They therefore do not imply
a condition of the form $\delta\ge n/2$. Substituting a Hamiltonian
theorem requiring a degree bound proportional to $n$ would leave a
quantifier gap.

\paragraph{A family showing why thresholds below one fail.}
For $1\le a\le b$ and $b\ge2$, the complete bipartite graph $K_{a,b}$
can be disconnected only by deleting all vertices of one side, since
any surviving vertices on both sides induce a connected graph.
Deleting the side of size $a$ produces $b$ isolated vertices and
ratio $a/b$. If additional vertices of the other side are deleted,
the ratio becomes $(a+r)/(b-r)\ge a/b$ whenever at least two
vertices remain. Deleting the other side first gives ratios at least
$b/a\ge a/b$ when that deletion is a qualifying cut. Thus its
toughness is exactly $a/b$.

A bipartite cycle visits equally many vertices in the two parts, so
$K_{a,b}$ is Hamiltonian exactly when $a=b\ge2$. In particular
$K_{r,r+1}$ is non-Hamiltonian and has toughness $r/(r+1)$, tending
to one from below. For every $t<1$, some member is $t$-tough and
non-Hamiltonian. This excludes all those thresholds, but the sequence
is bounded and gives no disproof of the existence of a larger threshold.

\paragraph{An exact obstruction above one.}
Use the Petersen graph on the ten two-element subsets of
$\{0,1,2,3,4\}$, with adjacency defined by disjointness. The accompanying
script enumerates all $2^{10}=1024$ vertex-deletion sets. For each
one producing at least two components it computes $|S|/c(G-S)$ as
an exact rational number. The minimum is
\[
                             \tau(P)=\frac43.
\]
With the subsets in lexicographic order, the cut with indices
$0,1,2,3$ attains the minimum and leaves three components. These
indices correspond to the four subsets containing zero; the remaining
six two-subsets of a four-element set form three disjoint edges.

The script independently proves non-Hamiltonicity of this particular
finite graph by exhaustive path dynamic programming. A state $(S,v)$
records a simple path beginning at vertex zero, using exactly $S$,
and ending at $v$. The initial state is $(\{0\},0)$; a state extends
to every unused neighbour of $v$. By induction this records every
simple path from zero. No state using all ten vertices ends at a
neighbour of zero, so no Hamiltonian cycle exists. This supplies a
finite certificate by complete enumeration, not a random search.
Consequently a universal threshold, if one exists, must be strictly
greater than $4/3$. No claim is made that this elementary obstruction
is the strongest known lower bound.

\paragraph{Why the two extremal inequalities still fall short.}
Toughness simultaneously controls every cut, whereas a proposed cycle
must choose compatible entrances and exits through all relevant cuts.
Even a bound on the independence number and a bound on connectivity
do not reconstruct those choices. Conversely, attaching many isolated
gadgets to a fixed separator makes Hamiltonicity fail but also makes
the toughness small: if a separator of size $s$ exposes $r$ pieces,
then toughness is at most $s/r$. Such a construction needs a proof
that every separating set scales with the number of pieces before it
can yield unbounded toughness.

\paragraph{Verification record and remaining gap.}
The exact Petersen cut ratios and Hamiltonian-path calculation are in
\texttt{checks\_graphs.py} and \texttt{results\_graphs.json} under
\path{_Local/Erdos_problems/UnsolvedMath/attempt_audit_2026_10_01/number_theory/}.
The argument excludes thresholds through $4/3$ and proves several
necessary consequences of any larger threshold. It proves neither
Hamiltonicity at a finite universal value nor unbounded toughness of
non-Hamiltonian graphs. That uniform alternative is the remaining
problem. \textbf{Outcome: UNRESOLVED.}

\subsection{Sources}

{\small

\begin{itemize}

\item[\hypertarget{source:1415:1}{S1}] ULAM.AI, \emph{UnsolvedMath}, supplied \texttt{problems.json}, record 1415 (GRAPH-028), statement and background. The embedded literature-review date is 2026-08-17; that is the archive’s review date, not a fresh certification. Dataset landing page: \url{https://huggingface.co/datasets/ulamai/UnsolvedMath}.

\item[\hypertarget{source:1415:2}{S2}] Hamiltonicity of 1-tough (P2 union kP1)-free graphs, Discrete Mathematics 348 (2025). \url{https://doi.org/10.1016/j.disc.2023.114082}. Bibliographic information imported from the supplied record.

\item[\hypertarget{source:1415:3}{S3}] Hamiltonicity of 1-tough (P2 union kP1)-free graphs, arXiv:2303.09741. \url{https://arxiv.org/abs/2303.09741}. Bibliographic information imported from the supplied record.


\item[E1] \emph{Hamiltonicity of $1$-tough
$(P_2\cup kP_1)$-free graphs}, arXiv:2303.09741.
\url{https://arxiv.org/abs/2303.09741}. Primary manuscript record
inspected 1 October 2026. Its forbidden-induced-subgraph hypothesis
is not part of the unrestricted toughness conjecture.

\end{itemize}

}

\label{end:1415}
\end{document}
